\documentclass[11pt,a4paper]{article}
\usepackage[T1]{fontenc}
\usepackage{mathptmx}
\usepackage[margin=25mm]{geometry}
\usepackage{booktabs}
\usepackage{array}
\usepackage{amsmath}
\usepackage{microtype}
\usepackage{titlesec}
\usepackage[hidelinks]{hyperref}
\usepackage{fancyhdr}
\titleformat{\section}{\large\bfseries}{\thesection}{0.6em}{}
\titleformat{\subsection}{\bfseries}{\thesubsection}{0.6em}{}
\setlength{\parskip}{0.5em}
\setlength{\parindent}{0pt}
\pagestyle{fancy}\fancyhf{}
\renewcommand{\headrulewidth}{0.4pt}
\newcommand{\inr}[1]{INR~#1}

\fancyhead[L]{RescueSwarm --- Track D: Perception}
\fancyhead[R]{\thepage}
\begin{document}

\begin{center}
{\LARGE\bfseries Track D --- Perception}\\[3mm]
{\large RescueSwarm: work package and funding request}\\[2mm]
NIDAR 2026--27 $\cdot$ Track 1 $\cdot$ Mission 1\\[1mm]
\textbf{Ask: \inr{76,245}} $\cdot$ 1--2 students $\cdot$ Detector, tiling, geotagging, calibration, dataset
\end{center}
\vspace{2mm}\hrule\vspace{4mm}


This track owns the 250 marks available for detection. It also owns the
programme's most honest weakness: with the ground-truth apparatus and the field
dataset both deferred, detection recall is a \emph{modelled} quantity and there
is currently no funded route to measuring it.

\section{Scope and ownership}

This track owns \textbf{detector, tiling, geotagging, calibration, dataset}. It is one of four delivery tracks defined in
\texttt{docs/development-plan.md}~\S1.3; the integrated system argument lives in
the master proposal, \texttt{docs/proposal/rescueswarm-proposal.pdf}, which this
document does not restate.

\section{What the detector actually receives}

The sensor is a 12.3\,MP module at $4056\times3040$ on a \textbf{1.55\,\textmu m}
pixel pitch behind a 6\,mm lens. The sizing chapter assumed a
\textbf{1.82\,\textmu m} pitch --- the same pixel \emph{count} at a different
pixel \emph{size}, which means resolution on target is not what that chapter
states.

At the 40\,m search altitude, taking a supine adult as 1.7\,m by 0.5\,m:

\begin{center}
\begin{tabular}{@{}lrrl@{}}
\toprule
\textbf{Stage} & \textbf{GSD} & \textbf{Target} & \textbf{COCO class} \\
 & (cm/px) & (px$^2$) & \\
\midrule
Sensor, native & 1.03 & 13\,600 & large \\
After $2\times$ downsample & 2.07 & 3\,400 & medium \\
Sizing chapter, 60\,m, $2\times$ & 3.65 & 640 & \textbf{small} \\
\bottomrule
\end{tabular}
\end{center}

The last row matters more than it looks. Published recall figures report average
precision \emph{by size class}, and small-object AP is routinely far below
medium on the same model. Quoting the 47-pixel figure implicitly promises
small-object performance the programme does not need to accept.

\section{The inference budget, and where it does not close}

Tiling is 640\,px tiles at 20\,\% overlap. At $2\times$ downsample that is 12
tiles per frame, and at 2\,Hz, \textbf{24 inferences per second} --- inside the
accelerator's capability.

\textbf{The temporal budget does not close.} Requirement SYS-46 asks for at
least 12 looks per target per pass so that multi-frame fusion has enough
independent observations. At 40\,m and 2\,Hz the target is in frame for
\textbf{7.9 looks} --- a 34\,\% shortfall. Meeting 12 looks needs
\[
r \geq \frac{n\,v}{D} = 3.06~\text{Hz},
\]
which is 37 inferences per second rather than 24. The trade is real and
unresolved: raise the capture rate and pay about half again in inference load,
fly slower, fly higher, or relax the look count. It should be decided
deliberately rather than discovered.

The overlap geometry does hold: at 128\,px of shared margin against an 82\,px
target, no survivor can be cut by a tile boundary without appearing whole in the
neighbouring tile.

\section{Geolocation: calibrate before fusing}

The error budget makes an ordering argument that this track should follow
literally. Random terms fall as $1/\sqrt{N}$ with frame count; systematic terms
do not fall at all. Boresight calibration and a ground-plane height fix
therefore matter \emph{more} than additional frames, and fusion is worth having
but not worth over-investing in --- it moves the total from 0.75\,m to 0.66\,m,
while calibration moves it from 4.57\,m to 3.88\,m before RTK is even applied.

Calibrate, then fuse. Not the other way round.

\section{The honest position on recall}

Detection recall is modelled and will remain modelled unless the deferred lines
are restored. The ground-truth apparatus would have given detections something
to be measured against; the field dataset would have given the detector Indian
terrain to be trained on. Both were deferred at team direction to reach the
requested figure.

This track should state that plainly in every review rather than presenting
modelled recall as though it were measured. If any single deferred line is
restored, this is the one with the largest effect on a scored outcome.

\section{Funding request}

Every figure below is generated from
\texttt{docs/proposal/figures/track\_budget.py}, which partitions the master
competition budget. The four track asks plus the shared pool reconcile to the
master figure exactly; that reconciliation is asserted in code and fails the
build if it ever stops holding. \textbf{K} marks a line held at professional
grade on purpose.


\begin{center}
\begin{tabular}{@{}p{88mm}cr@{}}
\toprule
\textbf{Item} & & \textbf{\inr{}} \\
\midrule
Camera + lens (1 x 11,100 x 3 ac) & \textbf{K} & 33,300 \\
AI accelerator (1 x 11,000 x 3 ac) & \textbf{K} & 33,000 \\
\midrule
Subtotal & & 66,300 \\
Customs duty and freight & & 0 \\
GST & & 0 \\
Contingency at 15\% & & 9,945 \\
\midrule
\textbf{TRACK D ASK} & & \textbf{76,245} \\
\bottomrule
\end{tabular}
\end{center}


\textbf{Deferred or institutionally held, and therefore not requested:}
Ground-truth apparatus, Indian SAR field dataset, Spare camera + lens, Training compute. These remain capabilities the track requires; they are simply not being
bought. Where a deferral carries a technical consequence it is stated in the
risk table below rather than left implicit.

\section{Deliverables by phase}

\begin{center}
\begin{tabular}{@{}llp{88mm}@{}}
\toprule
\textbf{Phase} & \textbf{Dates} & \textbf{This track delivers} \\
\midrule
\textbf{P1} & 24 Aug -- 13 Sep & Camera and accelerator ordered. Survey-altitude position stated with its evidence. \\
\textbf{P2} & 14 Sep -- 4 Oct & Detector running on the accelerator at the budgeted rate. Boresight calibration procedure defined. \\
\textbf{P3} & 5 Oct -- 18 Oct & First airborne imagery at the flown altitude. \\
\textbf{P5} & 9 Nov -- 22 Nov & Detection and geotagging in the loop, end to end. \\
\bottomrule
\end{tabular}
\end{center}

\section{Risks owned by this track}

\begin{center}
\begin{tabular}{@{}p{50mm}lp{70mm}@{}}
\toprule
\textbf{Risk} & \textbf{Sev.} & \textbf{Mitigation} \\
\midrule
Recall is unmeasured & High & Ground truth and dataset both deferred. Mitigation is honesty in review plus opportunistic collection during flight tests; the real fix costs money and is a restoration decision. \\
Fusion look count fails SYS-46 at the flown altitude & High & 7.9 looks against 12 required. Needs a decision on capture rate, groundspeed or altitude. Shared with Track C. \\
Fine-tuning at the wrong target scale & Medium & Altitude error changes pixels on target, so the detector should be fine-tuned across a band around 82\,px rather than trained at a single scale. \\
\bottomrule
\end{tabular}
\end{center}

\section{Interfaces to other tracks}

\begin{center}
\begin{tabular}{@{}lp{110mm}@{}}
\toprule
\textbf{With} & \textbf{Dependency} \\
\midrule
Track A & Depends on a stable camera mounting plane. Boresight drift is a systematic error that no amount of frame fusion removes. \\
Track B & Depends on the accelerator, the camera interface, and on time and attitude being accurate enough that a geotag means something. \\
Track C & Supplies detections; consumes aircraft state. Shares the unresolved survey-altitude decision. \\
\bottomrule
\end{tabular}
\end{center}

Interfaces are where student programmes fail, because each side assumes the
other owns the gap. Every row above is a two-way commitment and should be
recorded as an interface control document at P1.


\vfill
\hrule\vspace{2mm}
{\small Generated by \texttt{tools/proposal/build\_track\_proposals.py} from
\texttt{docs/proposal/figures/track\_budget.py}. Financial figures are not
maintained in this document and must not be edited here: change the master
budget and regenerate. Technical claims cite the model or document that
produced them.}
\end{document}
